%
% Introduction
%
\chapter{Introduction} \label{sec:introduction}
Seminal work in reinforcement learning (RL) has demonstrated that autonomous agents can achieve superhuman performance in complex sequential decision-making tasks.
Landmark successes in games such as Atari, Go, StarCraft II, and Dota 2 have shown that RL agents are capable of learning elaborate strategies directly through interaction with an environment~\cite{silver2017alphazero,vinyals2019starcraft,berner2019dota}.
These achievements have established reinforcement learning as a powerful approach to solving problems in which optimal decisions depend on long-term consequences rather than immediate rewards~\cite{silver2017alphazero}.

Among this research it is typically assumed that the environment remains stationary throughout both training and deployment~\cite{sutton2018rl,ditzler2015nonstationary}.
Under this crucial assumption the states, transition dynamics, available actions, and reward structure do not change over time, allowing an agent to gradually converge towards an effective policy.
In practice, however, many real-world environments are inherently non-stationary: Autonomous systems operating in finance, robotics, recommendation systems, or online games must continually adapt to changing conditions without forgetting previously acquired knowledge~\cite{ditzler2015nonstationary,khetarpal2022crl}.

Competitive multiplayer games provide an intuitive example of non-stationary environments.
Unlike traditional board games, live-service games receive frequent balance patches that modify game mechanics, adjust character abilities, introduce new content, or remove existing strategies~\cite{chitayat2023gameupdates}.
These updates intentionally alter the strategic landscape to maintain competitive balance and player engagement~\cite{chitayat2023gameupdates}.
Consequently, reinforcement learning agents trained on one version of the game may experience degraded performance when evaluated on subsequent versions.
Alternatively, training a reinforcement learning agent from scratch after every update is computationally expensive and might discard valuable knowledge acquired from previous versions of the game.

One promising direction for addressing this problem is continual reinforcement learning (CRL)~\cite{khetarpal2022crl,wang2024cl}.
Continual learning aims to enable an agent to acquire new knowledge while retaining previously learned capabilities, allowing it to adapt to a sequence of tasks over time~\cite{wang2024cl}.
Continual learning methods seek to mitigate catastrophic forgetting, the tendency of neural networks to overwrite previously learned information when trained on new tasks, while facilitating positive transfer between related tasks~\cite{wang2024cl,mccloskey1989catastforget}.
If successful, continual reinforcement learning could significantly reduce retraining costs while maintaining competitive performance as environments evolve.

This thesis aims to investigate continual reinforcement learning in the context of competitive Pokémon, also known as VGC.
VGC is an official Esports competition in which players construct teams of Pokémon under a set of rules that changes periodically through game updates, downloadable content, and seasonal rulesets~\cite{angliss2025vgc,vgc-rules-regulations}.
In addition to introducing new Pokémon, moves, and mechanics, these updates substantially reshape the competitive metagame by altering which strategies are viable~\cite{angliss2025vgc}.
Explicitly defined rulesets and task boundaries make VGC a compelling domain for studying continual learning in a non-stationary environment.

Compared to other commonly used reinforcement learning benchmarks, VGC presents several additional challenges.
Battles involve imperfect information, stochastic outcomes, an extremely large combinatorial state space, and long-term strategic planning~\cite{angliss2025vgc}.
Changes between rulesets often preserve some strategic knowledge while rendering other aspects obsolete, creating opportunities for knowledge transfer, but also increasing the risk of catastrophic forgetting.

The objective of this thesis is to investigate whether continual reinforcement learning methods can facilitate the effective transfer of reinforcement learning agents across successive VGC rulesets.
We investigate the impact of changing rulesets and evaluate whether knowledge learned in earlier game versions can be transferred to newer versions without requiring a complete retraining.
We compare continual learning approaches with conventional retraining in terms of competitive performance, generalization over unseen scenarios, and computational training cost.
To address these objectives, we define and investigate the following problem statement and research questions:

\paragraph{Problem Statement:} How can we train reinforcement learning agents on competitive games that receive frequent updates, with each update possibly invalidating previously learned strategies?

\begin{enumerate}[leftmargin=3em]
    \item[\textbf{RQ1:}] What effect does non-stationarity have on the performance and generalization of reinforcement learning agents in VGC?
    \item[\textbf{RQ2:}] Can continual reinforcement learning methods facilitate effective knowledge transfer across game versions?
    \item[\textbf{RQ3:}] How do continual reinforcement learning-based agents compare to fully retrained agents in terms of performance, generalization, and training cost across game updates?
\end{enumerate}

By investigating continual reinforcement learning in a realistic, evolving, competitive game environment, this thesis aims to contribute to a broader understanding of how reinforcement learning systems can remain effective in non-stationary domains.

The rest of this thesis is structured as follows:
We will first introduce relevant background work in Chapter \ref{sec:background}.
In Chapter \ref{sec:preliminary} we will perform preliminary experiments on the effects of non-stationarity and the transferability of our environment and baseline agents.
Based on these results, we will outline our proposed solutions in Chapter \ref{sec:methodology}, and perform our main experiments in Chapter \ref{sec:experiments}.
We will discuss the results of these experiments in Chapter \ref{sec:discussion}, and end by outlining our final conclusions and proposing directions for future research in Chapter \ref{sec:conclusion}.
